\chapter{Probability spaces}
\label{ch:spaces}

\begin{question}
What exactly is ``a probability''? What are we allowed to ask about a random experiment, and what rules must the answers obey?
\end{question}

\begin{whatfor}
\begin{itemize}
  \item Builds the three pieces of every model in this book: outcomes $\Omega$, events $\F$, probability $P$.
  \item Shows how a growing collection of events records growing information (a first look at filtrations).
  \item Proves the rules of probability we use everywhere, continuity of $P$, and the Borel--Cantelli lemma.
\end{itemize}
\skipline{fs:spaces}
\end{whatfor}

\begin{prereq}
Nothing beyond sets and convergent series (see ``Before you start'').
\end{prereq}

%=====================================================================
\section{Outcomes and events}

Run a random experiment. Exactly one result comes out.

\begin{definition}[Sample space, outcome]
The \dfn{sample space} $\Omega$ is the set of all possible results of the experiment. Each element $\omega\in\Omega$ is an \dfn{outcome}. When the experiment runs, one outcome $\omega$ is produced.
\end{definition}

\begin{example}[Three sample spaces]
These three will keep coming back.
\begin{itemize}
  \item Roll a die: $\Omega=\{1,2,3,4,5,6\}$. You roll a 4, so $\omega=4$.
  \item Toss a coin twice: $\Omega=\{HH,HT,TH,TT\}$. Heads then tails gives $\omega=HT$.
  \item Pick a point at random in $[0,1]$: $\Omega=[0,1]$, infinitely many outcomes. You get, say, $\omega=0.37$.
\end{itemize}
\end{example}

\begin{definition}[Event]
An \dfn{event} is a subset $A\subseteq\Omega$. We say $A$ \emph{happened} if the outcome lies in it, $\omega\in A$, and that $A$ did not happen if $\omega\notin A$.
\end{definition}

\begin{example}[Did the event happen?]
Die, $A=\{2,4,6\}$ (``even'').
\begin{itemize}
  \item Roll 4: $\omega=4\in A$, so $A$ happened.
  \item Roll 3: $\omega=3\notin A$, so $A$ did not happen; its complement $A^c=\{1,3,5\}$ (``odd'') happened instead.
\end{itemize}
$\Omega$ always happens; the empty set $\varnothing$ never happens.
\end{example}

\begin{pic}
\begin{tikzpicture}[scale=0.9]
  \draw[thick,defc,rounded corners] (0,0) rectangle (6,3);
  \node[defc] at (0.4,2.65) {$\Omega$};
  \draw[thick,thmc,fill=thmc!10] (2.2,1.5) ellipse (1.5 and 0.9);
  \node[thmc] at (1.3,1.9) {$A$};
  \fill (2.6,1.3) circle (2pt) node[right] {$\omega$};
  \fill (5,0.7) circle (2pt) node[right] {$\omega'$};
\end{tikzpicture}
\picnote{If the experiment produces $\omega$, the event $A$ happened. If it produces $\omega'$, $A$ did not happen and $A^c$ did.}
\end{pic}

%=====================================================================
\section{\texorpdfstring{$\sigma$}{Sigma}-fields: the events we may ask about}

We want to give each event a probability. For mathematical reasons we cannot always give one to \emph{every} subset of $\Omega$, so we fix in advance a collection of events we are allowed to ask about. The collection must be closed under the obvious logical operations: if we can ask ``did $A$ happen?'', we can ask ``did $A$ not happen?''.

\begin{definition}[$\sigma$-field]
\label{def:sigma-field}
A \dfn{$\sigma$-field} on $\Omega$ is a collection $\F$ of subsets of $\Omega$ such that
\begin{enumerate}[label=(\roman*)]
  \item $\varnothing\in\F$;
  \item if $A\in\F$, then $A^c=\Omega\setminus A\in\F$;
  \item if $A_1,A_2,\ldots\in\F$, then $\bigcup_{i=1}^\infty A_i\in\F$.
\end{enumerate}
\end{definition}

Two consequences come free.
\begin{itemize}
  \item $\Omega=\varnothing^c\in\F$.
  \item Intersections: $A\cap B=(A^c\cup B^c)^c$ is a complement of a union of complements, so it is in $\F$. The same holds for countably many sets.
\end{itemize}

\begin{example}[A $\sigma$-field with four events]
\label{ex:die-even-odd}
Die, $\F=\{\varnothing,\{2,4,6\},\{1,3,5\},\Omega\}$: you only learn ``even or odd''.
\begin{steps}
  \item $\varnothing\in\F$.
  \item Complements: $\{2,4,6\}^c=\{1,3,5\}\in\F$ and $\varnothing^c=\Omega\in\F$.
  \item Unions: $\{2,4,6\}\cup\{1,3,5\}=\Omega\in\F$; any union with $\varnothing$ or $\Omega$ is again in $\F$.
\end{steps}
So $\F$ is a $\sigma$-field. On the other hand $\{\varnothing,\{6\},\Omega\}$ is \emph{not} one: the complement $\{6\}^c=\{1,2,3,4,5\}$ is missing.
\end{example}

\begin{pic}
\begin{tikzpicture}
  \foreach \k in {1,...,6} {
    \pgfmathparse{mod(\k,2)==1 ? "thmc" : "defc"}\edef\c{\pgfmathresult}
    \draw[thick,draw=\c,fill=\c!15] (\k*1.1,0) circle (0.4) node {\k};
  }
  \node[thmc,anchor=west] at (8,0.25) {\small odd $=\{1,3,5\}$};
  \node[defc,anchor=west] at (8,-0.25) {\small even $=\{2,4,6\}$};
\end{tikzpicture}
\picnote{The $\sigma$-field of \cref{ex:die-even-odd} cuts $\Omega$ into two blocks. Every event in it is a union of whole blocks: none, one, or both.}
\end{pic}

This picture is the general pattern on a finite $\Omega$. A $\sigma$-field cuts $\Omega$ into disjoint non-empty blocks, called \dfn{atoms}, and its events are exactly the unions of atoms.

\begin{result}[Counting events]
\label{res:count}
If a $\sigma$-field on a finite $\Omega$ has $k$ atoms, it has exactly $2^k$ events. The smallest $\sigma$-field is $\{\varnothing,\Omega\}$ (one atom, $2$ events); the largest is the collection of all subsets of $\Omega$ ($2^{|\Omega|}$ events).
\end{result}
\begin{proof}
Each event is a union of atoms, and each atom is either in it or out of it: $2^k$ choices. Different choices give different sets, because the atoms are disjoint and non-empty.
\end{proof}

For the die: even/odd has $2$ atoms and $2^2=4$ events; all subsets has $6$ atoms and $2^6=64$ events.

%=====================================================================
\section{Information that grows: a first look at filtrations}
\label{sec:info}

Toss a coin infinitely many times, so $\Omega$ is the set of all infinite sequences $\omega=(\omega_1,\omega_2,\ldots)$ of $H$ and $T$. Write $A_H$ for ``the first toss is $H$'', $A_{HT}$ for ``the first two tosses are $H$ then $T$'', and so on.

After $k$ tosses you have seen $\omega_1,\dots,\omega_k$. Some events you can now decide for sure (``$A_H$ happened''); others you cannot yet (``$A_{HH}$'' needs the second toss when $k=1$).

\begin{definition}[Information at time $k$]
$\F_k$ is the collection of all events that can be decided, for sure, from the first $k$ tosses.
\end{definition}

\begin{example}[The first three levels]
\label{ex:coin-filtration}
Each level is a $\sigma$-field, cut into finer atoms as $k$ grows.
\begin{itemize}
  \item $\F_0=\{\varnothing,\Omega\}$: nothing seen yet.
  \item $\F_1=\{\varnothing,A_H,A_T,\Omega\}$: two atoms, $4$ events.
  \item $\F_2$: atoms $A_{HH},A_{HT},A_{TH},A_{TT}$, so $2^4=16$ events, for example $A_{HH}\cup A_{HT}=A_H$.
\end{itemize}
$A_{HH}$ is in $\F_2$ but not in $\F_1$: after seeing $H$ first you still cannot tell $HH$ from $HT$.
\end{example}

\begin{pic}
\begin{tikzpicture}[scale=0.85]
  \def\w{8}
  \draw[thick] (0,2.4) rectangle (\w,3); \node at (\w/2,2.7) {$\Omega$}; \node[left] at (0,2.7) {$\F_0$};
  \draw[thick] (0,1.2) rectangle (\w/2,1.8); \draw[thick] (\w/2,1.2) rectangle (\w,1.8);
  \node at (\w/4,1.5) {$A_H$}; \node at (3*\w/4,1.5) {$A_T$}; \node[left] at (0,1.5) {$\F_1$};
  \foreach \i/\l in {0/HH,1/HT,2/TH,3/TT} {\draw[thick] (\i*\w/4,0) rectangle ({(\i+1)*\w/4},0.6); \node at ({(\i+0.5)*\w/4},0.3) {$A_{\l}$};}
  \node[left] at (0,0.3) {$\F_2$};
  \draw[-{Latex},defc,thick] (\w+0.4,3)--(\w+0.4,0) node[midway,right,align=left] {\small more\\\small tosses\\\small seen};
\end{tikzpicture}
\picnote{The same $\Omega$, cut finer and finer. Each new toss splits every atom in two, so information only grows: $\F_0\subseteq\F_1\subseteq\F_2\subseteq\cdots$.}
\end{pic}

\begin{definition}[Filtration]
An increasing sequence of $\sigma$-fields $\F_0\subseteq\F_1\subseteq\F_2\subseteq\cdots$ is called a \dfn{filtration} (also a flow of $\sigma$-fields). $\F_k$ represents all the information available by time $k$.
\end{definition}

Filtrations return in Part~1, where they become the main tool.

%=====================================================================
\section{Borel sets}

On $\Omega=\R$ we want to ask whether a number lands in an interval, and anything built from intervals.

\begin{definition}[Generated $\sigma$-field; Borel sets]
Given any collection $\mathcal C$ of subsets of $\Omega$, the $\sigma$-field \dfn{generated by} $\mathcal C$, written $\sigma(\mathcal C)$, is the smallest $\sigma$-field containing every set in $\mathcal C$. The $\sigma$-field generated by all intervals of $\R$ is written $\B(\R)$; its members are the \dfn{Borel sets}.
\end{definition}

Every interval, every single point $\{x\}=\bigcap_n(x-\frac1n,x+\frac1n)$, every countable set, and every finite union of intervals is a Borel set. For this book: \emph{any set of numbers you can write down is a Borel set}.

%=====================================================================
\section{Probability measure and probability space}

\begin{definition}[Probability measure, probability space]
\label{def:prob}
Let $\F$ be a $\sigma$-field on $\Omega$. A \dfn{probability measure} is a function $P:\F\to[0,1]$ with
\begin{enumerate}[label=(\roman*)]
  \item $P(\Omega)=1$;
  \item if $A_1,A_2,\ldots\in\F$ are pairwise disjoint ($A_i\cap A_j=\varnothing$ for $i\ne j$), then $P\big(\bigcup_{i=1}^\infty A_i\big)=\sum_{i=1}^\infty P(A_i)$.
\end{enumerate}
The triple $(\Omega,\F,P)$ is a \dfn{probability space}. The sets in $\F$ are the events of the model.
\end{definition}

\begin{intuition}
One line each: $\Omega$ is what can happen; $\F$ is what you are allowed to ask; $P$ is how likely each answer is.
\end{intuition}

\begin{example}[One fair coin, written out completely]
$\Omega=\{H,T\}$, $\F=\{\varnothing,\{H\},\{T\},\{H,T\}\}$, and $P(\varnothing)=0$, $P(\{H\})=P(\{T\})=\frac12$, $P(\{H,T\})=1$.
\end{example}

\begin{result}[Rules of probability]
\label{res:rules}
For events $A,B,A_1,A_2,\ldots$:
\begin{enumerate}[label=(\alph*)]
  \item $P(A^c)=1-P(A)$, and $P(\varnothing)=0$;
  \item if $A\subseteq B$ then $P(A)\le P(B)$;
  \item $P(A\cup B)=P(A)+P(B)-P(A\cap B)$;
  \item \dfn{union bound}: $P\big(\bigcup_i A_i\big)\le\sum_i P(A_i)$.
\end{enumerate}
\end{result}
\begin{proof}
Each rule cuts a set into disjoint pieces and uses rule (ii).
\begin{steps}
  \item $A$ and $A^c$ are disjoint with union $\Omega$, so $P(A)+P(A^c)=1$. With $A=\Omega$: $P(\varnothing)=0$.
  \item $B=A\cup(B\setminus A)$, disjoint, so $P(B)=P(A)+P(B\setminus A)\ge P(A)$.
  \item $A\cup B=A\cup(B\setminus A)$ and $B=(A\cap B)\cup(B\setminus A)$, both disjoint. Subtract the two equations.
  \item Make the sets disjoint: $C_1=A_1$, $C_i=A_i\setminus(A_1\cup\dots\cup A_{i-1})$. Then $\bigcup C_i=\bigcup A_i$ and $C_i\subseteq A_i$, so $P(\bigcup A_i)=\sum P(C_i)\le\sum P(A_i)$. \qedhere
\end{steps}
\end{proof}

The union bound is loose by exactly the overlaps it double-counts. Die, $A=\{1,2\}$, $B=\{2,3\}$: $P(A\cup B)=\frac12\le\frac13+\frac13$.

%=====================================================================
\section{Almost surely}

\begin{definition}[Almost surely]
An event $A$ happens \dfn{almost surely} (a.s.) if $P(A)=1$; equivalently $P(A^c)=0$.
\end{definition}

``Almost surely'' is weaker than ``surely''. Surely means $A=\Omega$: no outcome lies outside $A$. Almost surely allows outcomes outside $A$, as long as together they have probability $0$.

\begin{example}[Possible, but of probability zero]
Pick a uniform point in $[0,1]$, so $P([a,b])=b-a$.
\begin{itemize}
  \item $A=\{\omega\ne0.5\}$. Then $P(\{0.5\})=P([0.5,0.5])=0$, so $P(A)=1$: $A$ happens a.s. But $\omega=0.5$ is still a possible outcome, so not surely.
  \item A coin tossed forever: $P(\text{first $n$ tosses all $H$})=(\frac12)^n\to0$, so ``at least one $T$'' happens a.s., although $HHH\ldots$ is a valid outcome.
\end{itemize}
On a finite $\Omega$ where every outcome has positive probability, a.s.\ and surely coincide; that is why die examples never show the difference.
\end{example}

%=====================================================================
\section{Continuity of probability}

\begin{theorem}[Continuity of $P$]
\label{thm:continuity}
Two monotone limits may be swapped with $P$.
\begin{enumerate}[label=(\alph*)]
  \item From below: if $A_1\subseteq A_2\subseteq\cdots$, then $P\big(\bigcup_n A_n\big)=\lim_{n\to\infty}P(A_n)$.
  \item From above: if $A_1\supseteq A_2\supseteq\cdots$, then $P\big(\bigcap_n A_n\big)=\lim_{n\to\infty}P(A_n)$.
\end{enumerate}
\end{theorem}
\begin{proof}
Cut the growing sets into rings.
\begin{steps}
  \item (a) Let $C_1=A_1$ and $C_n=A_n\setminus A_{n-1}$. The rings are disjoint, $A_n=C_1\cup\dots\cup C_n$, and $\bigcup_n A_n=\bigcup_n C_n$. So
  \[
    P\Big(\bigcup_n A_n\Big)=\sum_{n=1}^\infty P(C_n)=\lim_{n\to\infty}\sum_{k=1}^n P(C_k)=\lim_{n\to\infty}P(A_n).
  \]
  \item (b) If $A_n$ shrink, their complements grow, and $\big(\bigcap_n A_n\big)^c=\bigcup_n A_n^c$. Apply (a) to $A_n^c$ and use $P(A^c)=1-P(A)$. \qedhere
\end{steps}
\end{proof}

Before swapping $\lim$ and $P$, check the nesting: the sets must only grow or only shrink.

%=====================================================================
\section{The Borel--Cantelli lemma}

Take events $A_1,A_2,\ldots$ and ask: which outcomes lie in infinitely many of them? Let
\[
  B_n=\bigcup_{i=n}^\infty A_i=A_n\cup A_{n+1}\cup\cdots,\qquad E=\bigcap_{n=1}^\infty B_n .
\]
\begin{itemize}
  \item $\omega\in B_n$ means: $\omega$ lies in at least one of $A_n,A_{n+1},\ldots$. As $n$ grows, $B_n$ drops the early sets, so $B_1\supseteq B_2\supseteq\cdots$.
  \item $\omega\in E$ means: for every $n$, $\omega$ lies in some $A_i$ with $i\ge n$. That is, $\omega$ lies in \emph{infinitely many} $A_i$.
\end{itemize}
$E$ is written $\limsup A_n$, or ``$A_n$ infinitely often''.

\begin{theorem}[Borel--Cantelli lemma]
\label{thm:bc}
If $\sum_{i=1}^\infty P(A_i)<\infty$, then $P(E)=0$. That is, almost surely only finitely many of the $A_i$ happen.
\end{theorem}
\begin{proof}
Three ingredients, in this order.
\begin{steps}
  \item $B_n$ shrink to $E$, so by continuity from above (\cref{thm:continuity}), $P(E)=\lim_n P(B_n)$.
  \item By the union bound, $P(B_n)\le\sum_{i=n}^\infty P(A_i)$.
  \item This is the tail of a convergent series, so it tends to $0$. Hence
  \[
    0\le P(E)=\lim_{n\to\infty}P(B_n)\le\lim_{n\to\infty}\sum_{i=n}^\infty P(A_i)=0. \qedhere
  \]
\end{steps}
\end{proof}

\begin{example}[Shrinking intervals]
\label{ex:bc}
Uniform point in $[0,1]$, $A_i=[0,1/i^2]$.
\begin{steps}
  \item $\sum_i P(A_i)=\sum_i 1/i^2=\pi^2/6\approx1.645<\infty$, so the lemma gives $P(E)=0$.
  \item Directly: the sets shrink, so $B_n=A_n=[0,1/n^2]$ and $E=\bigcap_n[0,1/n^2]=\{0\}$, which has probability $0$.
  \item Meaning: an outcome $\omega>0$ lies in $A_i$ only while $1/i^2\ge\omega$. For $\omega=0.01$ that is $i\le10$; after that cut-off, $A_i$ never happens again.
\end{steps}
\end{example}

\begin{pic}
\begin{tikzpicture}[xscale=10,yscale=0.42]
  \draw[-{Latex}] (0,0)--(1.08,0) node[right] {$\omega$};
  \foreach \x in {0,0.25,0.5,1} \draw (\x,0.15)--(\x,-0.15) node[below] {\scriptsize $\x$};
  \foreach \i in {1,...,6} {
    \pgfmathsetmacro{\r}{1/(\i*\i)}
    \draw[line width=3pt,defc!80] (0,\i) -- (\r,\i);
    \node[left] at (0,\i) {\scriptsize $A_{\i}$};
  }
  \draw[dashed,trapc] (0.01,0)--(0.01,7);
  \node[trapc,right] at (0.012,6.8) {\scriptsize $\omega=0.01$ is in $A_1,\dots,A_{10}$ only};
\end{tikzpicture}
\picnote{The sets $A_i=[0,1/i^2]$ shrink towards $\{0\}$. Every $\omega>0$ falls out after finitely many steps; only $\omega=0$ stays in all of them, and $P(\{0\})=0$.}
\end{pic}

\begin{remark}
The lemma goes one way only. A divergent sum tells you nothing on its own: $A_i=[0,1/i]$ has $\sum 1/i=\infty$, yet again $E=\{0\}$ and $P(E)=0$.
\end{remark}

%=====================================================================
\formulasheet{fs:spaces}

\begin{mastercard}[label=card:atoms]{Events of a finite $\sigma$-field}
A finite $\sigma$-field is fixed by its atoms $G_1,\dots,G_k$: its smallest non-empty events, disjoint and covering $\Omega$. Every event is a union of atoms:
\[
  \boxed{\;\F=\Big\{\textstyle\bigcup_{i\in I}G_i\Big\},\qquad |\F|=2^k,\qquad P\Big(\textstyle\bigcup_{i\in I}G_i\Big)=\sum_{i\in I}P(G_i).\;}
\]
\begin{plugrows}
  \plugrow{Is a given collection a $\sigma$-field?}{\pstep{Box: a $\sigma$-field is exactly all unions $\bigcup_{i\in I}G_i$ of its atoms (disjoint, covering $\Omega$).}\pstep{From the list: take its smallest non-empty sets as candidate atoms.}\pstep{Check they are disjoint and cover $\Omega$.}\pstep{Check every union of them is in the list (complements first: they fail most often).}}{yes exactly when the list is all unions of some disjoint sets that cover $\Omega$}{finite $\Omega$}{die, list $\{\varnothing,\{6\},\Omega\}$: smallest non-empty sets $\{6\},\Omega$ are not disjoint; the complement $\{6\}^c=\{1,\dots,5\}$ is missing, so \textbf{no}}
  \plugrow{How many events?}{\pstep{Box: $|\F|=2^k$.}\pstep{From the atoms: $k$ = their number.}\pstep{Substitute $k$.}}{$2^k$ events (each atom is in or out of the union)}{always (finite)}{die, even/odd: $k=2$, $2^2=4$ events; all subsets: $k=6$, $2^6=64$}
  \plugrow{Information after $n$ coin tosses, $\F_n$}{\pstep{Box: $|\F|=2^k$.}\pstep{From ``what is seen'': the atoms are the sets of sequences sharing the first $n$ tosses, so $k=2^n$.}\pstep{Substitute: $|\F_n|=2^{k}=2^{2^n}$.}\pstep{Each atom of $\F_n$ is a union of atoms of $\F_{n+1}$, so $\F_n\subseteq\F_{n+1}$.}}{$|\F_n|=2^{2^n}$; $\F_0\subseteq\F_1\subseteq\cdots$ (a filtration)}{tosses observed one by one}{$\F_1$: $k=2^1=2$ atoms $A_H,A_T$, $2^2=4$ events; $\F_2$: $k=4$, $2^4=16$ events}
  \plugrow{Smallest / largest $\sigma$-field}{\pstep{Box: $|\F|=2^k$, events = unions of atoms.}\pstep{Smallest: one atom, $G_1=\Omega$, so $k=1$.}\pstep{Largest: every outcome is an atom, so $k=|\Omega|$.}}{$\{\varnothing,\Omega\}$ ($2^1=2$ events) / all subsets ($2^{|\Omega|}$ events)}{finite $\Omega$}{die: $2^1=2$ and $2^6=64$ events}
  \plugrow{Probability of an event in $\F$}{\pstep{Box: $P\big(\bigcup_{i\in I}G_i\big)=\sum_{i\in I}P(G_i)$.}\pstep{From the event: write it as a union of atoms, which gives $I$.}\pstep{Substitute: add the probabilities of those atoms.}}{$P(\text{event})=\sum$ of its atoms' probabilities}{atoms disjoint}{fair die, atoms $\{1,3,5\},\{2\},\{4\},\{6\}$: $\{1,2,3,5\}=\{1,3,5\}\cup\{2\}$, so $P=\frac36+\frac16=\frac12+\frac16=\frac23$}
\end{plugrows}
\end{mastercard}

\begin{mastercard}[label=card:incl-excl]{Probability of a union}
For any events $A_1,\dots,A_n$:
\[
  \boxed{\begin{gathered}P\Big(\bigcup_{i=1}^nA_i\Big)=\sum_iP(A_i)-\sum_{i<j}P(A_i\cap A_j)+\sum_{i<j<k}P(A_i\cap A_j\cap A_k)-\cdots\;\le\;\sum_iP(A_i).\end{gathered}}
\]
\begin{plugrows}
  \plugrow{Disjoint events}{\pstep{Box: $P(\bigcup A_i)=\sum P(A_i)-\sum_{i<j}P(A_i\cap A_j)+\cdots$.}\pstep{From ``disjoint'': every intersection is $\varnothing$, so every $P(A_i\cap A_j)=0$, and so on.}\pstep{Substitute: only the first sum is left.}}{$P(\bigcup A_i)=\sum P(A_i)$}{pairwise disjoint (also for countably many)}{die: $P(\{1\}\cup\{2\})=\frac16+\frac16-0=\frac13$}
  \plugrow{Complement}{\pstep{Box with $n=2$: $P(A_1\cup A_2)=P(A_1)+P(A_2)-P(A_1\cap A_2)$.}\pstep{From the complement: $A_1=A$, $A_2=A^c$, so $A_1\cup A_2=\Omega$ and $A_1\cap A_2=\varnothing$.}\pstep{Substitute: $1=P(A)+P(A^c)-0$.}\pstep{With $A=\Omega$: $P(\varnothing)=1-1=0$.}}{$P(A^c)=1-P(A)$; $P(\varnothing)=0$}{always}{$P(\text{not }6)=1-P(6)=1-\frac16=\frac56$}
  \plugrow{Two events}{\pstep{Box with $n=2$: the first sum has two terms, the second sum one term, no more terms.}}{$P(A\cup B)=P(A)+P(B)-P(A\cap B)$}{always}{die, $A=\{1,2\}$, $B=\{2,3\}$, $A\cap B=\{2\}$: $\frac26+\frac26-\frac16=\frac36=\frac12$}
  \plugrow{Upper bound (union bound)}{\pstep{Box, right-hand end: $P(\bigcup A_i)\le\sum_iP(A_i)$.}\pstep{The terms after the first sum only remove overlaps counted twice.}}{$P(\bigcup A_i)\le\sum P(A_i)$}{always; exact iff overlaps have probability 0}{same $A,B$: $P(A\cup B)=\frac12\le P(A)+P(B)=\frac13+\frac13=\frac23$}
  \plugrow{Subset}{\pstep{Disjoint case (row 1): $P(A\cup C)=P(A)+P(C)$ for disjoint $A,C$.}\pstep{From $A\subseteq B$: $B=A\cup(B\setminus A)$, and $A$, $B\setminus A$ are disjoint.}\pstep{Substitute: $P(B)=P(A)+P(B\setminus A)\ge P(A)$.}}{$A\subseteq B\Rightarrow P(A)\le P(B)$}{always}{$\{6\}\subseteq\{\text{even}\}$: $P(\{6\})=\frac16\le P(\text{even})=\frac36=\frac12$}
\end{plugrows}
\end{mastercard}

\begin{mastercard}[label=card:limits]{Limits of events}
If the sets only grow ($A_n\uparrow$) or only shrink ($A_n\downarrow$), probability and limit swap:
\[
  \boxed{\begin{gathered}P\Big(\lim_{n\to\infty}A_n\Big)=\lim_{n\to\infty}P(A_n),\\ \lim A_n=\bigcup_nA_n\ (\uparrow)\ \text{ or }\ \bigcap_nA_n\ (\downarrow).\end{gathered}}
\]
\begin{plugrows}
  \plugrow{Growing sets}{\pstep{Box: $P(\lim A_n)=\lim P(A_n)$.}\pstep{From ``growing'': $\lim A_n=\bigcup_nA_n$.}\pstep{Substitute.}}{$P(\bigcup A_n)=\lim P(A_n)$}{$A_1\subseteq A_2\subseteq\cdots$}{uniform on $[0,1]$, $A_n=[0,1-\frac1n]$: $\bigcup A_n=[0,1)$, so $P([0,1))=\lim(1-\frac1n)=1-0=1$}
  \plugrow{Shrinking sets}{\pstep{Box: $P(\lim A_n)=\lim P(A_n)$.}\pstep{From ``shrinking'': $\lim A_n=\bigcap_nA_n$.}\pstep{Substitute.}}{$P(\bigcap A_n)=\lim P(A_n)$}{$A_1\supseteq A_2\supseteq\cdots$}{$A_n$ = ``first $n$ tosses all $H$'', $P(A_n)=2^{-n}$: $P(\text{all }H)=\lim2^{-n}=0$}
  \plugrow{Infinitely often (Borel--Cantelli)}{\pstep{Box, shrinking case: $P(\bigcap_nB_n)=\lim P(B_n)$.}\pstep{From ``infinitely often'': $B_n=\bigcup_{i\ge n}A_i$ shrinks to $\{A_n\text{ i.o.}\}$.}\pstep{Union bound (\cref{card:incl-excl}, row 4): $P(B_n)\le\sum_{i\ge n}P(A_i)$.}\pstep{Substitute: $P(A_n\text{ i.o.})=\lim P(B_n)\le\lim\sum_{i\ge n}P(A_i)=0$, the tail of a convergent sum.}}{if $\sum_iP(A_i)<\infty$: $P(A_n\text{ i.o.})=0$; a.s.\ only finitely many $A_n$ happen}{convergent sum; a divergent sum proves nothing}{$A_i=[0,1/i^2]$, $P(A_i)=\frac1{i^2}$: $\sum\frac1{i^2}=\frac{\pi^2}6\approx1.645<\infty$, so $P(A_n\text{ i.o.})=0$ (indeed the limit set is $\{0\}$)}
  \plugrow{Almost surely}{\pstep{Definition: $A$ holds a.s.\ when $P(A)=1$, i.e.\ $P(A^c)=0$ (\cref{card:incl-excl}, row 2).}\pstep{From the problem: find the exceptional set $A^c$.}\pstep{Show $P(A^c)=0$, often as a shrinking limit (row 2).}}{$A$ holds a.s.\ $\iff P(A)=1$}{a.s.\ allows outcomes in $A^c$; on a finite $\Omega$ with all $P(\omega)>0$ a.s.\ = surely}{uniform on $[0,1]$, $A=\{\omega\ne0.5\}$: $A^c=\{0.5\}=\bigcap_n[0.5-\frac1n,0.5+\frac1n]$, $P=\lim\frac2n=0$, so $\omega\ne0.5$ a.s.}
\end{plugrows}
\end{mastercard}
